\chapter{Derived Hardware Numerology}
\label{app:derived_hardware_numerology}

The Applied Vacuum Engineering (AVE) framework operates under the strictest formulation of first-principles physics. There are precisely three arbitrary calibration inputs ($\ell_{node}$, $\alpha$, and $G$) defining the entire universal manifold. Every subsequent physical constraint—particularly those governing the limits of topological continuous-wave frameworks—must be analytically traceable back to these parameters without empirical curve-fitting.

This appendix derives the fundamental numerology appearing throughout Volume 4, verifying their exact theoretical limits against the continuous LC tensor solver (used in Chapters 2-18) and the declarative compilation traces (Chapter 19-20).

\section{The Vacuum Impedance ($Z_0$)}

The base characteristic impedance of the vacuum lattice limits all perfectly matched Klopfenstein tapers and topological routing structures ($\Gamma=0$).
\begin{equation}
    Z_0 = \sqrt{\frac{\mu_0}{\epsilon_0}} \approx 376.730313668 \;\Omega
\end{equation}
This serves as the exact physical bounding resistance for Chiral Circulator shunts preventing spatial reflections.

\section{The Dielectric Rupture Voltage ($V_{snap}$)}

The Axiom 4 node fracture limit separates pure geometry from material particle dynamics. This establishes the Reverse Breakdown saturation boundary for topological continuous limiters.
\begin{equation}
    V_{snap} = \frac{m_e c^2}{e} \approx 510{,}998.950 \text{ V}
\end{equation}
As proven during the declarative compilation tests (Chapter 20), this value operates purely as a structural spatial constraint replacing the specific chemical junction breakdown limits in conventional CMOS hardware.

\section{The Kinetic Yield Point ($V_{yield}$)}

Before reaching the absolute fracture boundary ($V_{snap}$), the lattice initiates continuous non-linear yielding—the primary mechanism required for Topological Logic states (Chapter 14) and Tensor Plate refractive gradients (Chapter 15).
This yield occurs intrinsically scaled by the Fine-Structure Constant ($\alpha$), derived dynamically in Volume~8.
\begin{equation}
    V_{yield} = \sqrt{\alpha} \cdot V_{snap} \approx \sqrt{\frac{1}{137.036}} \cdot 510{,}998 \text{ V} \approx 43{,}652.7 \text{ V}
\end{equation}

\section{Geometric Strain Yield ($\phi_{yield}$)}

Converting the localized Kinetic Yield Point into continuous structural strain dictates exactly how Non-Volatile Soliton Kinks (Chapter 8) are permanently packed.
In purely geometric terms, topological packing maps strictly to the tetrahedral $sp^3$ limit yielding:
\begin{equation}
    \phi_{yield} = \arcsin\left(\frac{\sqrt{3}}{2}\right) = \frac{\pi}{3} \text{ radians} \approx 1.047 \text{ rad}
\end{equation}
When the geometry is driven to hold $\phi > \pi/3$, the internal node connections cross the Axiom 4 limit, physically generating the stable topological knot.

\section{FDTD Numerical Damping Factor}

The FDTD Topo-Kinematic solver employs a standard numeric \texttt{sponge\_damping = 0.8} profile within boundary zones explicitly to absorb arbitrary calculation artifact reflections. 

\textit{Important Note:} This $0.8$ scalar is a \emph{numerical stability artefact} necessary for finite discrete array mathematics inside the simulation engine. It is not an Axiomatic property of the universe, but rather limits artificial grid boundary explosions during FDTD evaluation. It is documented here explicitly to denote it is not a ``magic number'' mapping true physics, but an analytical constraint of the \texttt{ave.core} solver framework.

\section{Corrugation Loading Factor ($\gamma$)}
\label{app:gamma_corrugation}
\textit{Derived in: Chapter~\ref{ch:dielectric_delay_lines} (Slow-Wave Structures).}

For a sinusoidal corrugation of pitch $\Lambda$ and depth $h$ etched into a waveguide of nominal width $w \gg h$, the true surface path length per unit axial length exceeds the nominal by the arc-length factor:
\begin{equation}
    \gamma = \frac{1}{\Lambda}\int_0^\Lambda \sqrt{1 + \left(\frac{dh\,\pi}{\Lambda}\cos\frac{2\pi x}{\Lambda}\right)^2}\,dx \approx 1 + \frac{\pi^2 h^2}{2\Lambda^2}
\end{equation}
The approximation is the leading term of the Taylor expansion valid for $h \ll \Lambda$. Both $L'_{\rm eff}$ and $C'_{\rm eff}$ increase by $\gamma$, so impedance $Z = \sqrt{L'/C'}$ is exactly preserved while phase velocity $v_{ph} = 1/\sqrt{L'C'}$ drops to $c_0/\gamma$. The excess propagation delay inserted by a corrugated segment of nominal length $\ell$ is:
\begin{equation}
    \Delta\tau = \frac{(\gamma-1)\,\ell}{c_0} \approx \frac{\pi^2 h^2 \ell}{2\Lambda^2 c_0}
\end{equation}
\textit{Axiom trace:} Axiom~1 (LC constitutive equations) → Telegraphist equations → $v_{ph} = 1/\sqrt{L'C'}$ → corrugation path-length geometry.

\section{Klopfenstein Ripple Parameter ($A$)}
\label{app:klopf_A}
\textit{Derived in: Chapter~\ref{ch:axiomatic_transducers} (Axiomatic Transducers).}

For an impedance taper from $Z_{min}$ to $Z_{max}$ with a passband equi-ripple reflectance target $\Gamma_{max}$, the Dolph-Chebyshev parameter $A$ is:
\begin{equation}
    A = \cosh^{-1}\!\left(\frac{\ln(Z_{max}/Z_{min})}{2\,\Gamma_{max}}\right)
\end{equation}
For the canonical VCA-to-RF interface ($Z_{min} = 50\,\Omega$, $Z_{max} = 376.73\,\Omega$, $\Gamma_{max} = 0.01$):
\begin{equation}
    A = \cosh^{-1}\!\left(\frac{\ln(7.535)}{0.02}\right) = \cosh^{-1}(100.95) \approx 5.31
\end{equation}
$A$ is not a free parameter — it is uniquely determined by the impedance ratio (set by Axiom~1) and the engineering reflectance specification. \\
\textit{Axiom trace:} Axiom~3 ($\Gamma$ at boundaries) → distributed reflection Fourier integral → Chebyshev equi-ripple optimality.

\section{Minimum Klopfenstein Taper Length at $f_{CC}$ ($L_{min}$)}
\label{app:L_klopf}
\textit{Derived in: Chapter~\ref{ch:axiomatic_transducers}.}

The minimum physical taper length for passband cutoff at the carrier coherence frequency $f_{CC}$ is:
\begin{equation}
    L_{min} = \frac{A\,c_0}{2\pi f_{CC}\sqrt{\kappa_{topo}}} = \frac{A\,\lambda_{op}}{2\pi}
\end{equation}
For the SOI substrate at the 3\,nm node: $f_{CC} = 1.832\,\text{THz}$, $\kappa_{topo} = 3.9$, giving $\lambda_{op} = c_0/(f_{CC}\sqrt{\kappa_{topo}}) \approx 83\,\mu\text{m}$, and therefore:
\begin{equation}
    L_{min} = \frac{5.31 \times 83\,\mu\text{m}}{2\pi} \approx 70\,\mu\text{m}
\end{equation}
Note: the free-space calculation gives $\lambda_{op} \approx 163\,\mu\text{m}$ ($L_{min} \approx 138\,\mu\text{m}$); the SOI dielectric reduces this by $\sqrt{\kappa_{topo}} \approx 1.97$. The correct engineering value accounts for the substrate $\kappa_{topo}$. \\
\textit{Axiom trace:} Axiom~1 ($v_{ph} = c_0/\sqrt{\kappa}$) → $f_{CC} = c_0/(2L\sqrt{\kappa})$ → $\lambda_{op}$ → taper length from $A$.

\section{Sine-Gordon Kink Write Energy ($U_{kink}$)}
\label{app:U_kink}
\textit{Derived in: Chapter~\ref{ch:static_soliton_kinks} (Static Soliton Kinks).}

The exact rest energy stored in a single static sine-Gordon kink (topological charge $Q = +1$) is evaluated by integrating the Hamiltonian density over the kink profile $\phi_{kink}(x) = 4\arctan(\exp((x-x_0)/\lambda_{kink}))$:
\begin{align}
    U_{kink} &= \int_{-\infty}^{+\infty} \left[\frac{c_0^2}{2}\!\left(\frac{d\phi}{dx}\right)^2 + \omega_0^2(1-\cos\phi)\right]dx \notag\\
    &= 8\,c_0\,\omega_0
\end{align}
The first-integral identity ($(d\phi/dx)^2 = (2\omega_0/c_0)^2\sin^2(\phi/2)$) equates both terms, doubling the kinetic integral. In physical VCA units:
\begin{equation}
    U_{kink} = 8\sqrt{\mu_{visc}\,\kappa_{topo}}\;V_{snap}
\end{equation}
This is the minimum energy required to nucleate a persistent topological memory bit. Once written, the kink requires $0.0\,\text{W}$ to maintain. \\
\textit{Axiom trace:} Axiom~4 → sine-Gordon restoring force → static kink ODE → first integral → Hamiltonian.

\section{Topo-Kinematic Write Threshold Voltage ($\Phi_{pack}\,V_{snap}$)}
\label{app:write_threshold}
\textit{Derived in: Chapter~\ref{ch:static_soliton_kinks}.}

The minimum write amplitude required to nucleate a stable soliton kink (raise the phase field past the $\phi = \pi$ potential barrier in the sine-Gordon double-well) is:
\begin{equation}
    V_{write} = \Phi_{pack}\,V_{snap} = \frac{\pi\sqrt{2}}{6} \times \frac{m_e c^2}{e} \approx 0.7405 \times 511\,\text{kV} \approx 378\,\text{kV}
\end{equation}
$\Phi_{pack} = \pi\sqrt{2}/6$ is the face-centered cubic (FCC) packing fraction. It enters here because Axiom~1 selects FCC as the ground-state geometry of the vacuum LC network ($K = 2G$ bulk-to-shear modulus ratio uniquely chooses FCC among Bravais lattices). The scaling $\Phi_{pack}\,V_{snap}$ is the fraction of the saturation amplitude at which the FCC lattice's topological void structure permits a stable $2\pi$ phase defect without global lattice reorganization — i.e., the first amplitude at which a kink can be localized rather than radiating away. \\
\textit{Axiom trace:} Axiom~1 (FCC ground state, $K=2G$) → $\Phi_{pack} = \pi\sqrt{2}/6$ → Axiom~4 ($V_{snap} = m_e c^2/e$) → product sets kink nucleation threshold.

%%%
\section{Viscous Drag Power Generation ($P_{drag}$)}
\label{app:viscous_drag_power}
\textit{Derived in: Volume 4, Chapter 15.}

In the Topo-Kinematic framework, power dissipation does not occur via Ohmic electron scattering ($I^2R$). Instead, it arises purely from the viscous drag of the continuous-wave carrier propagating through an imperfect crystalline dielectric (the substrate). Because velocity in the Topo-Kinematic identity equates to the time derivative of the phase field ($\vec{v} \leftrightarrow \dot{\phi}$), the power loss is strictly proportional to the excitation frequency $\omega$, not the squared amplitude of a rest-mass current.

The analytic loss profile integrated over the bulk topological volume $V$ is:
\begin{equation}
    P_{drag} = \omega_{cc}\, \mu_{visc}\, \kappa_{topo}\, \tan\delta \iiint_V |S_{carrier}|^2 dV
\end{equation}
Where $\omega_{cc} = 2\pi f_{CC}$ is the continuous-wave carrier frequency, $\mu_{visc} \leftrightarrow$ vacuum inductance, $\kappa_{topo} \leftrightarrow$ substrate permittivity relative to vacuum, and $\tan\delta$ is the loss tangent of the substrate material. 

Evaluating this analytic integral for a $10\,\text{mm} \times 10\,\text{mm}$ active substrate layer with a focal volume utilization ratio of $10^{-6}$ and an underlying SOI loss tangent of $\tan\delta \approx 10^{-4}$ at $f_{CC} = 1.832\,\text{THz}$ yields exactly:
\begin{equation}
    P_{drag} \approx 19.8\text{ W}
\end{equation}
This confirms that computing at THz scales without melting requires transitioning to high-purity ($\tan\delta \ll 1$) topological substrates where heat is a linear $f^1$ scaling rather than the quadratic $CV^2f$ CMOS catastrophe. \\
\textit{Axiom trace:} Axiom~1 ($\vec{v} \leftrightarrow i$, $\mu_{visc}$ metric) → dielectric loss tangent integration over VCA manifold.

%%%
\section{Soliton Kink Node Density ($\rho_{kink}$)}
\label{app:kink_density}
\textit{Derived in: Chapter~\ref{ch:axiomatic_processing_unit_capstone}.}

The peak spatial density of non-volatile topological memory is established by the 2D cross-sectional area footprint of a single sine-Gordon soliton kink. As derived in Chapter~\ref{ch:static_soliton_kinks}, the kink has a characteristic full-width of $2\lambda_{kink}$, where $\lambda_{kink} = c_0/\omega_0$ is the structural attenuation length.

The theoretical maximum areal density prior to catastrophic inter-knot structural interference is:
\begin{equation}
    \rho_{kink} = \frac{1}{(2\lambda_{kink})^2}
\end{equation}
Submitting the exact $\omega_0$ rest-frame eigenfrequency of the Topo-Kinematic knot vacuum state into the JAX characterization suite yields an effective $\lambda_{kink} \approx 2.4 \times 10^{-14}\,\text{m}$. Taking the inverse square provides the peak packing density limit:
\begin{equation}
    \rho_{kink} \approx 4.34 \times 10^{20} \text{ knots/mm}^2
\end{equation}
This analytical capacity exceeds contemporary Flash/SRAM density constraints by over ten orders of magnitude because topological state storage operates at the fundamental vacuum-lattice cutoff rather than the lithographic electron-confinement limits of silicon. \\
\textit{Axiom trace:} Axiom~4 ($U_{kink}$ geometry) → sine-Gordon kink width $2\lambda_{kink}$ → 2D spatial footprint density.

\section{Adler Phase-Degeneracy Tolerances ($\lambda / 8$)}
\label{app:adler_headroom}
\textit{Derived in: Volume 4, Chapter 15.}

In synchronization manifolds like the 2-Beam Interference Coupler, signal recombination is successful if and only if the temporal offset remains bounded by the Adler synchronization envelope.

The required headroom preventing stochastic phase slip diverges exactly at $1/8^{\text{th}}$ of a spatial wavelength:
\begin{equation}
    \Delta x_{\max} = \frac{\lambda_{op}}{8} = \frac{\pi}{4k}
\end{equation}
Operating outside this $\lambda / 8$ phase boundary forces nodes into the nonlinear coupling regime where deterministic signal recombination becomes structurally suppressed. This spatial invariant dictates physical trace length match requirements on the routing board regardless of substrate type. \\
\textit{Axiom trace:} Axiom~3 (Adler-locking dynamics) → $\pi/2$ phase margin relative displacement → $\lambda/8$ focal depth limit.

\section{The Trace-Reversed Poisson Ratio ($\nu_{vac}$)}
\label{app:trace_reversed_poisson}
\textit{Derived in: Volume~1, Universal Operators.}

The scalar Refractive Index lensing operator (Op 19) leverages the trace-reversed geometry of the symmetric spatial coordinate system. The dimensionless property locking metric strain to the refractive gradient is exactly:
\begin{equation}
    \nu_{vac} = \frac{2}{7} \approx 0.2857
\end{equation}
Validating that $n(r) = 1 + \nu_{vac} \varepsilon_{11}(r)$ completely satisfies the exact Schwartzschild orbital deflection bounds without requiring arbitrary relativistic fitting parameters. \\
\textit{Axiom trace:} Axiom~1 (Machian symmetric elasticity) → Cauchy boundary conditions → $\nu_{vac} = 2/7$.

\section{The Effective Coordination Number ($z_0$)}
\label{app:z0_coordination}
\textit{Derived in: Volume~1, Chapter~1; Appendix~\ref{app:full_derivation_chain}, Layer~3.}

The effective coordination number of the chiral vacuum lattice emerges as the physical root of the Feng-Thorpe-Garboczi EMT quadratic. Setting the packing fraction $p^* = 8\pi\alpha$ equal to the trace-reversal condition $K/G = 2$:
\begin{equation}
    p^* = \frac{10\,z_0 - 12}{z_0(z_0 + 2)} = 8\pi\alpha
\end{equation}
Rearranging: $\;8\pi\alpha\, z_0^2 + (16\pi\alpha - 10)\,z_0 + 12 = 0$. The physical root is:
\begin{equation}
    z_0 = \frac{10 - 16\pi\alpha + \sqrt{(10-16\pi\alpha)^2 - 4 \cdot 8\pi\alpha \cdot 12}}{2 \cdot 8\pi\alpha}
    \approx \boxed{51.25}
\end{equation}
The second root ($z_0 \approx 1.28$) is unphysical---3D rigidity requires $z_0 \geq 6$. The rigidity threshold at this coordination is $p_G = 6/z_0 \approx 0.117$, confirming the vacuum operates $56.7\%$ above the fluid-solid boundary. \\
\textit{Axiom trace:} Axiom~4 ($\alpha \equiv p_c/8\pi$) → EMT $K/G = 2$ → quadratic solve → $z_0 \approx 51.25$.

\section{The Macroscopic Avalanche Exponent ($n_{3D}$)}
\label{app:avalanche_n_3d}
\textit{Derived in: Volume~3, Chapter~\ref{ch:kolmogorov_spectral_cutoff}.}

The nonlinear transfer of energy through a fully developed turbulent cascade amplifies according to the universal avalanche factor $M = 1/\mathcal{S}^2 = 1/(1-r^2)$, which directly derives from Axiom 4 power conservation (an effective pure 1D exponent of $n=2$, directly mapping Tabletop Relativity $\gamma^2$ properties). For 3D isotropic shear flow, energy leakages into transverse lattice modes scale precisely by the vacuum Poisson ratio $\nu_{vac} = 2/7$.

The effective 3D exponent adjusts as:
\begin{equation}
    n_{3D} = 2 \left(1 - \frac{\nu_{\mathrm{vac}}}{3}\right) = 2\left(1 - \frac{2/7}{3}\right) = \frac{38}{21} \approx 1.8095
\end{equation}
\textit{Axiom trace:} Axiom 4 ($\mathcal{S}^2 + r^2 = 1$) → Avalanche $M = 1/\mathcal{S}^2 \rightarrow n=2$ → Axiom 3 ($\nu_{vac}=2/7$) → $n_{3D} = 38/21$.

\section{The Kolmogorov Constant ($C_K$)}
\label{app:kolmogorov_constant}
\textit{Derived in: Volume~3, Chapter~\ref{ch:kolmogorov_spectral_cutoff}.}

The universal Kolmogorov constant $C_K$, determining the magnitude of the inertial subrange turbulence spectrum, emerges strictly from the topological scattering properties of the continuous K4 mesh structure. For energy flux traversing any generalized crossing or junction vertex, the mesh matrix establishes pure lattice forward cascade efficiency $\eta$:
\begin{equation}
    \eta = 3_{ports} \times |S_{ij}|^2 = 3 \times \left(\frac{1}{2}\right)^2 = \frac{3}{4}
\end{equation}
Taking the inverse of this efficiency reveals the exact Kolmogorov constant:
\begin{equation}
    C_K = \frac{1}{\eta} = \frac{4}{3} \approx 1.333
\end{equation}
\textit{Axiom trace:} Axiom 3 (S-matrix for uniform isotropic lattice intersections) → cascade efficiency $\eta = 3/4$ $\rightarrow C_K = 4/3$.
